% ==================== 一、项目与设计目标 ====================
\section{项目与设计目标}

\subsection{项目选择与承接关系}

本项目是实验一「需求分析」所选的第 \textcircled{4} 个题目——\textbf{小型酒店预订管理系统}的延续。实验一交付的是需求规格说明：四类角色（顾客、前台、店长、平台地区负责人）、两台状态机（订单六态、房间六态）、按「房型 $\times$ 日期」逐晚记录的库存模型、以及并发抢房、改价与下单互斥、价格快照等跨模块约束，共 36 条需求（其中已确认 34 条，REQ-035、REQ-036 仍标「待确认」）。

实验二的定位不同：\textbf{不写任何业务代码}，只把需求翻译成可以照着实现的设计。因此本报告自始至终遵守一条原则——\textbf{凡是能落到表结构、接口签名、状态转换表上的，就写实；凡是只能停留在口头描述的，就当成没设计}。报告里出现的每一个类、每一张表、每一个界面状态，都能回答「它对应哪条需求」和「实验三照着它能不能写出来」。

\subsection{设计目标}

\begin{enumerate}[leftmargin=3.2em]
  \item \textbf{不超卖}。同一房型同一晚最后一份库存，被 100 个请求同时抢，只能卖出 1 间；多晚订单中任一晚不足则整单失败，不留半截占用。
  \item \textbf{状态一致}。订单与房间各有一台状态机，非法转换一律拒绝而不是「尽力修正」；支付、取消、超时、未到店四个入口改的是同一行数据，且都带前置状态条件，谁先改成谁生效。
  \item \textbf{关键操作可追溯}。改价、开单、入住、取消、审批、删除，每一次写入都留操作人、时间、变更前后值；日志不提供修改与删除入口。
  \item \textbf{每个模块都有界面}。这是实验二的交付约束，也是本设计在模块划分阶段实际调整过划分的依据（见第七章 U-01）。
\end{enumerate}

\subsection{设计范围}

\begin{table}[H]
\centering
\caption{本次设计覆盖的内容与深度}
\tabfont
\begin{tabularx}{\textwidth}{@{}L{3.0cm}L{3.6cm}X@{}}
\toprule
设计项 & 产物 & 深度 \\
\midrule
软件的主要功能 & 第 2 章 + 表 \ref{tab:trace} & 36 条需求逐条映射到模块，标明跨模块与不覆盖的需求 \\
架构设计 & 第 3 章 + 包图、组件图、部署图 & 分层依赖规则、模块边界、依赖方向、部署形态写死 \\
数据存储设计 & 第 4 章 + 12 张表 & 字段、主键、唯一约束、索引；幂等与库存的约束落在具体列上 \\
行为建模 & 第 5 章 + 用例图、状态图、活动图 & 状态转换表穷举，非法转换给出处理方式 \\
交互建模 & 第 5 章 + 4 张时序图 & 抢房、支付回调、改价互斥、超时释放与对账 \\
UI 设计 & 第 6 章 + 16 屏界面 & 四端各屏的元素、状态与对应需求 \\
\bottomrule
\end{tabularx}
\end{table}

\subsection{明确不做的部分}

实验一已把会员、积分、优惠券、OTA 对接、在线开票划到范围之外（REQ-019），本设计同样不涉及。此外还有三处设计层面主动放弃的能力，写在这里以免读者以为漏了：

\begin{itemize}
  \item \textbf{不做微服务拆分}。本系统日均订单在百量级、门店十余间房，拆成多个服务只会把「扣库存 + 落单」这一个事务拆成两段，破坏不超卖的前提。模块边界按领域切，但部署为一个应用进程，见 3.5 决策一。
  \item \textbf{不做多活与异地容灾}。REQ-017 只要求营业时间可用性 99.5\%、故障 2 小时内恢复，主从 + 备份 + binlog 回放已经足够，多活会引入跨地域写冲突，与「库存只在一处串行扣减」直接矛盾。
  \item \textbf{不做动态定价}。REQ-023 明确「房价不与剩余房间数量挂钩」，因此系统里没有收益管理算法，价格是人定的，系统只负责守住平台下发的区间。
\end{itemize}

\subsection{统一的数据基准}

报告里 16 张 UML 图和 16 屏界面用的是同一套数据，读者可以拿界面上的数字去反查类图与表结构。基准如下（快照时刻 \textbf{2026-10-31 周六 13:47}）：

\begin{itemize}
  \item \textbf{门店}：湖滨银泰酒店，杭州西湖区，5 星，评分 4.6。
  \item \textbf{房型与楼层}：豪华双床房 24 间（1F、2F、6F），高级大床房 16 间（4F、5F），行政套房 8 间（3F），合计 48 间。
  \item \textbf{当日房态}（48 间）：在住 33、需打扫 7、已预定 3、维修中 2（408 水管、608 空调）、空闲 2、已退房 1，出租率 33/48 = 68.8\%。
  \item \textbf{当日价格}：高级大床房 ¥428/晚，豪华双床房 ¥498/晚，行政套房 ¥888/晚。
  \item \textbf{当日库存}（高级大床房）：总房量 16 − 线下保留 8 − 维修占用 1 = 线上可售上限 7，线上已售 3，剩余 4；次日 11-01 维修已结束，上限 8、已售 7、剩余 1，故「10-31 至 11-02 两晚」的可订量为 $\min(4,1)=1$ 间。
  \item \textbf{片区}：6 家门店、296 间房，当日营收 ¥113,721，加权出租率 72.7\%，平均房价（ADR）¥528.7。
\end{itemize}

\begin{quote}
\zihao{5}这套数字不是随手编的。它被三处交叉校验过：48 间的六种房态必须相加等于 48；48 间的价格 $\times$ 出租率必须能解释当日营收；片区 6 店的出租率按房量加权必须等于 72.7\%。界面上的每个数字改动都要过这三道校验，否则就是设计里出现了两套事实。
\end{quote}
